\BourbakiExercisesSection

\begin{enumerate}
\def\labelenumi{\arabic{enumi}.}
\tightlist
\item
  设 K 为一个域， \(K'\) 为 K 的一个子域，而 V 是 K 上的一个非零右向量空间。
\end{enumerate}

\begin{enumerate}
\def\labelenumi{(\alph{enumi})}
\item
  设 \(\mathbf{V}'\) 是一个 \(\mathbf{K}'\) -结构（在 \(\mathbf{V}\) 上）。证明 \(\mathbf{K}'\) 等于由下述 \(\mu \in \mathbf{K}\) 组成的集：对所有 \(x \in \mathbf{V}'\) , \(x\mu \in \mathbf{V}'\) 。
\item
  设 \(\Gamma'\) 为 K 的保持 K' 的元素不变的自同构群。证明：若不存在 K 的不属于 K' 且在所有自同构 \(\sigma \in \Gamma'\) 下不变的元素，则不存在 V 的不属于 V' 且在 V 的每个保持 V' 的所有元素不变的双射双态（相对于 K 的一个自同构）下不变的元素。
\item
  反之，设 E 为 V 的一个子集，G 为保持 E 的所有元素不变的 V 的双射双态群；对所有 \(u \in G\) ，设 \(\sigma_{u}\) 为 K 的相应自同构，并令 \(\Gamma\) 为 K 的自同构群，即 G 在同态 \(u \mapsto \sigma_{u}\) 下的像。假设 G 不包含 V 的不同于恒等自同构的自同构，并且不存在 V 的不属于 E 且在所有双态 \(u \in G\) 下不变的元素。证明：若 \(K'\) 是 K 的在所有自同构 \(\sigma \in \Gamma\) 下不变的元素所成的子域，则 E 是 V 上的一个 \(K'\) -结构。（首先证明 E 是 V 的一个加法子群且包含 V 的一组基；令 \(K''\) 为由下述元素 \(\mu \in K\) 组成的集：关系 \(x \in E\) 蕴含 \(x\mu \in E\) 。证明 \(K''\) 的元素在所有 \(\sigma \in \Gamma\) 下不变，并且 \(K''\) 是 K 的一个子域；由此推出 E 是 V 上的一个 \(K''\) -结构。）
\end{enumerate}

\begin{enumerate}
\def\labelenumi{\arabic{enumi}.}
\setcounter{enumi}{1}
\tightlist
\item
  设 K 为一个域， \(K'\) 为 K 的一个子域，V 是 K 上的一个右向量空间， \(V'\) 是 V 上的一个 \(K'\) -结构，而 \(R'\) 是对偶 \(V'*\) 在 V 的对偶 \(V*\) 中的典范像（第 4 节）。为使 \(R'\) 是 \(V*\) 的一个生成系，必要且充分条件是 V 是有限维的，或者 K 是一个有限维的右向量 \(K'\) -空间（利用 § 7 的习题 25）。
\end{enumerate}

¶ 3. 设 K 是一个域，L 是 K 的一个子域，令 \(K_{L}\) 表示把 K 视为 L 上的左向量空间而得到的域； \(E = \text{End}_{L}(K_{L})\) 典范地具有一个 (K, K)-双模结构。对偶 \((K_{L})^{*}\) 包含在 E 中，并且是 E 的一个 (K, L)-子双模。当 L 包含在 K 的中心内时，E 上的两个 L-向量空间结构——分别把 E 的两个 K-向量空间结构的标量域限制到 L 而得到——是相同的。

此后假设 \(\mathbf{L}\) 包含于 \(\mathbf{K}\) 的中心，且 \(\mathbf{K}_{\mathbf{L}}\) 的维数有限并等于 \(n\) 。

\begin{enumerate}
\def\labelenumi{(\alph{enumi})}
\item
  证明 \((\mathrm{K}_{\mathrm{L}})^{*}\) 在其 K 上的左向量空间结构下是 1 维的（用两种方式计算 \((\mathrm{K}_{\mathrm{L}})^{*}\) 在 L 上的维数）。
\item
  证明 \(\mathbf{E}\) 是一个 \(n\) 维左向量空间（在 \(\mathbf{K}\) 上）（同样的方法）。
\item
  设 F 是 K 上的一个有限维左向量空间，而 \(F_{[L]}\) 是 L 上的相应向量空间。若 \(u_{0} \neq 0\) 是 \(K_{L}\) 上的一个线性形式，证明映射 \(x' \mapsto u_{0} \circ x'\) （由 F 的对偶 \(F^{*}\) （把 F 视为 L 上的向量空间）映到对偶 \((\mathrm{F}_{[\mathrm{L}]})^{*}\) （即 \(F_{[L]}\) 的对偶）上）是双射的（利用 (a)）。
\end{enumerate}

¶ 4. 设 K 是其中心 Z 上有限秩的域，L 是 K 的包含 Z 的子域。

\begin{enumerate}
\def\labelenumi{(\alph{enumi})}
\item
  证明K关于L的左、右向量空间结构具有相同的维数。
\item
  证明习题 3 的性质 (a)、(b)、(c) 在这种情形也成立（若 \(u_0\) 是一个线性形式 \(\neq 0\) （在 \(\mathbf{L}_{\mathbf{Z}}\) 上），而 \(v_0\) 是一个线性形式 \(\neq 0\) （在 \(\mathbf{K}_{\mathbf{L}}\) 上），注意 \(\xi. (u_0 \circ v_0) = u_0 \circ (\xi v_0)\) （对 \(\xi \in \mathbf{K}\) ），并且 \(\xi. (u_0 \circ v_0)\) 遍历对偶 \((\mathbf{K}_{\mathbf{Z}})^*\) （当 \(\xi\) 遍历 \(\mathbf{K}\) 时），并利用习题 3 (c)）。
\end{enumerate}

\begin{enumerate}
\def\labelenumi{\arabic{enumi}.}
\setcounter{enumi}{4}
\tightlist
\item
  设 \(K_{0}\) 是域 K 的一个子域，使得 K 作为 \(K_{0}\) 上的右向量空间具有维数 2。设 E 是 K 上的一个左向量空间。
\end{enumerate}

\begin{enumerate}
\def\labelenumi{(\alph{enumi})}
\item
  设 \(E_{0}\) 是 E 的一个子集，它是 \(K_{0}\) 上的向量空间，并设 V 是 \(E_{0}\) 的最大向量子 K-空间；若 \(W_{0}\) 是一个向量子 \(K_{0}\) -空间（属于 \(E_{0}\) ，与 V 互补），证明由 \(W_{0}\) 生成的 E 的向量子 K-空间 W 满足 \(V \cap W = \{0\}\) （注意，若元素 \(\mu \in K\) 不属于 \(K_{0}\) ，则关系 \(\mu x \in E_{0}\) （对某个 \(x \in E_{0}\) ）蕴含 \(x \in V\) ）。
\item
  设 \(E_{0}^{\prime}\) 是 E 的另一个向量子 \(K_{0}\) -空间，而 \(V^{\prime}\) 是 \(E_{0}^{\prime}\) 的最大向量子 K-空间。为使存在 E 的一个把 \(E_{0}\) 映到 \(E_{0}^{\prime}\) 的 K-自同构，必要且充分条件是 V 和 \(V^{\prime}\) 关于 K 具有相同的维数，V 在 \(E_{0}\) 中与 \(V^{\prime}\) 在 \(E_{0}^{\prime}\) 中的余维数（关于 \(K_{0}\) ）相等，并且 \(E_{0}\) 与 \(E_{0}^{\prime}\) 在 E 中的余维数（关于 \(K_{0}\) ）相等（利用 (a)）。
\end{enumerate}

